% Options for packages loaded elsewhere
\PassOptionsToPackage{unicode}{hyperref}
\PassOptionsToPackage{hyphens}{url}
\PassOptionsToPackage{dvipsnames,svgnames,x11names}{xcolor}
%
\documentclass[
  10pt,
]{article}
\usepackage{amsmath,amssymb}
\usepackage{iftex}
\ifPDFTeX
  \usepackage[T1]{fontenc}
  \usepackage[utf8]{inputenc}
  \usepackage{textcomp} % provide euro and other symbols
\else % if luatex or xetex
  \usepackage{unicode-math} % this also loads fontspec
  \defaultfontfeatures{Scale=MatchLowercase}
  \defaultfontfeatures[\rmfamily]{Ligatures=TeX,Scale=1}
\fi
\usepackage{lmodern}
\ifPDFTeX\else
  % xetex/luatex font selection
  \setmainfont[]{DejaVu Serif}
  \setsansfont[]{DejaVu Sans}
  \setmonofont[]{DejaVu Sans Mono}
  \setmathfont[]{Latin Modern Math}
\fi
% Use upquote if available, for straight quotes in verbatim environments
\IfFileExists{upquote.sty}{\usepackage{upquote}}{}
\IfFileExists{microtype.sty}{% use microtype if available
  \usepackage[]{microtype}
  \UseMicrotypeSet[protrusion]{basicmath} % disable protrusion for tt fonts
}{}
\makeatletter
\@ifundefined{KOMAClassName}{% if non-KOMA class
  \IfFileExists{parskip.sty}{%
    \usepackage{parskip}
  }{% else
    \setlength{\parindent}{0pt}
    \setlength{\parskip}{6pt plus 2pt minus 1pt}}
}{% if KOMA class
  \KOMAoptions{parskip=half}}
\makeatother
\usepackage{xcolor}
\usepackage[margin=2.0cm]{geometry}
\usepackage{color}
\usepackage{fancyvrb}
\newcommand{\VerbBar}{|}
\newcommand{\VERB}{\Verb[commandchars=\\\{\}]}
\DefineVerbatimEnvironment{Highlighting}{Verbatim}{commandchars=\\\{\}}
% Add ',fontsize=\small' for more characters per line
\newenvironment{Shaded}{}{}
\newcommand{\AlertTok}[1]{\textcolor[rgb]{1.00,0.00,0.00}{\textbf{#1}}}
\newcommand{\AnnotationTok}[1]{\textcolor[rgb]{0.38,0.63,0.69}{\textbf{\textit{#1}}}}
\newcommand{\AttributeTok}[1]{\textcolor[rgb]{0.49,0.56,0.16}{#1}}
\newcommand{\BaseNTok}[1]{\textcolor[rgb]{0.25,0.63,0.44}{#1}}
\newcommand{\BuiltInTok}[1]{\textcolor[rgb]{0.00,0.50,0.00}{#1}}
\newcommand{\CharTok}[1]{\textcolor[rgb]{0.25,0.44,0.63}{#1}}
\newcommand{\CommentTok}[1]{\textcolor[rgb]{0.38,0.63,0.69}{\textit{#1}}}
\newcommand{\CommentVarTok}[1]{\textcolor[rgb]{0.38,0.63,0.69}{\textbf{\textit{#1}}}}
\newcommand{\ConstantTok}[1]{\textcolor[rgb]{0.53,0.00,0.00}{#1}}
\newcommand{\ControlFlowTok}[1]{\textcolor[rgb]{0.00,0.44,0.13}{\textbf{#1}}}
\newcommand{\DataTypeTok}[1]{\textcolor[rgb]{0.56,0.13,0.00}{#1}}
\newcommand{\DecValTok}[1]{\textcolor[rgb]{0.25,0.63,0.44}{#1}}
\newcommand{\DocumentationTok}[1]{\textcolor[rgb]{0.73,0.13,0.13}{\textit{#1}}}
\newcommand{\ErrorTok}[1]{\textcolor[rgb]{1.00,0.00,0.00}{\textbf{#1}}}
\newcommand{\ExtensionTok}[1]{#1}
\newcommand{\FloatTok}[1]{\textcolor[rgb]{0.25,0.63,0.44}{#1}}
\newcommand{\FunctionTok}[1]{\textcolor[rgb]{0.02,0.16,0.49}{#1}}
\newcommand{\ImportTok}[1]{\textcolor[rgb]{0.00,0.50,0.00}{\textbf{#1}}}
\newcommand{\InformationTok}[1]{\textcolor[rgb]{0.38,0.63,0.69}{\textbf{\textit{#1}}}}
\newcommand{\KeywordTok}[1]{\textcolor[rgb]{0.00,0.44,0.13}{\textbf{#1}}}
\newcommand{\NormalTok}[1]{#1}
\newcommand{\OperatorTok}[1]{\textcolor[rgb]{0.40,0.40,0.40}{#1}}
\newcommand{\OtherTok}[1]{\textcolor[rgb]{0.00,0.44,0.13}{#1}}
\newcommand{\PreprocessorTok}[1]{\textcolor[rgb]{0.74,0.48,0.00}{#1}}
\newcommand{\RegionMarkerTok}[1]{#1}
\newcommand{\SpecialCharTok}[1]{\textcolor[rgb]{0.25,0.44,0.63}{#1}}
\newcommand{\SpecialStringTok}[1]{\textcolor[rgb]{0.73,0.40,0.53}{#1}}
\newcommand{\StringTok}[1]{\textcolor[rgb]{0.25,0.44,0.63}{#1}}
\newcommand{\VariableTok}[1]{\textcolor[rgb]{0.10,0.09,0.49}{#1}}
\newcommand{\VerbatimStringTok}[1]{\textcolor[rgb]{0.25,0.44,0.63}{#1}}
\newcommand{\WarningTok}[1]{\textcolor[rgb]{0.38,0.63,0.69}{\textbf{\textit{#1}}}}
\usepackage{longtable,booktabs,array}
\usepackage{calc} % for calculating minipage widths
% Correct order of tables after \paragraph or \subparagraph
\usepackage{etoolbox}
\makeatletter
\patchcmd\longtable{\par}{\if@noskipsec\mbox{}\fi\par}{}{}
\makeatother
% Allow footnotes in longtable head/foot
\IfFileExists{footnotehyper.sty}{\usepackage{footnotehyper}}{\usepackage{footnote}}
\makesavenoteenv{longtable}
\setlength{\emergencystretch}{3em} % prevent overfull lines
\providecommand{\tightlist}{%
  \setlength{\itemsep}{0pt}\setlength{\parskip}{0pt}}
\setcounter{secnumdepth}{-\maxdimen} % remove section numbering
% doc_header.tex -- LaTeX preamble additions used by build_docs.sh for every
% Markdown -> LaTeX -> PDF conversion (supplements, README, submission files).
\usepackage{etoolbox}
% longtable may not be loaded by pandoc when a document has no table
\usepackage{longtable}
% bold vectors: \bm -> \symbf under unicode-math (XeLaTeX), bm package otherwise
\makeatletter
\@ifpackageloaded{unicode-math}{\providecommand{\bm}[1]{\symbf{#1}}}{\usepackage{bm}}
\makeatother
% smaller type in long tables so that wide audit tables fit the text width;
% inside a table a line break is allowed after an underscore of plain text
% (pandoc writes it as \_), so that long identifiers wrap inside their column
% rows are set a little apart, so that symbols with indices do not touch
\AtBeginEnvironment{longtable}{\footnotesize\renewcommand{\arraystretch}{1.45}%
  \let\sptunderscore\_\def\_{\sptunderscore\allowbreak}}
\setlength{\LTcapwidth}{\textwidth}
% allow line breaks after slashes and hyphens inside table cells
\usepackage{xurl}
% code blocks: break over-long lines instead of running into the margin.
% Plain ``` blocks are typeset by pandoc as {verbatim}, highlighted ones
% (```bash) as {Highlighting}; both are re-declared as breakable fancyvrb
% environments.  The continuation line is marked with a small hooked arrow.
\usepackage{fvextra}
\RecustomVerbatimEnvironment{verbatim}{Verbatim}{breaklines,breakanywhere,fontsize=\small}
\makeatletter
\@ifundefined{Highlighting}{}{%
  \RecustomVerbatimEnvironment{Highlighting}{Verbatim}{breaklines,breakanywhere,fontsize=\small,commandchars=\\\{\}}}
\makeatother
% inline code (pandoc writes it as \texttt{...}): never hyphenate -- an inserted
% hyphen could be read as part of a file name -- but allow a line break after
% an underscore, slash, comma, colon or equals sign, so that long paths and
% column lists wrap inside the text width and inside table columns.
% Inside a formula (e.g. the name of a notebook symbol in an equation) the same
% command must first leave math mode, which \text of amsmath does.
\ExplSyntaxOn
\tl_new:N \l__spt_code_tl
\cs_new_protected:Npn \__spt_texttt:n #1
  {
    \tl_set:Nn \l__spt_code_tl {#1}
    \regex_replace_all:nnN { ([\/\,\:\=]) } { \1 \c{allowbreak} } \l__spt_code_tl
    \regex_replace_all:nnN { \c{_} } { \c{_} \c{allowbreak} } \l__spt_code_tl
    \group_begin: \ttfamily \hyphenchar \font = -1 \scan_stop: \l__spt_code_tl \group_end:
  }
\RenewDocumentCommand \texttt { m }
  {
    \mode_if_math:TF { \text { \__spt_texttt:n {#1} } } { \__spt_texttt:n {#1} }
  }
\ExplSyntaxOff
% A figure is written in the Markdown source as an image paragraph followed by
% a caption paragraph ("**Figure S1.** ..."), so that the caption is also
% visible on GitHub.  An image that starts a paragraph is centered.
\usepackage{graphicx}% pandoc loads it only for documents that contain an image
\NewCommandCopy\sptincludegraphics\includegraphics
\RenewDocumentCommand\includegraphics{O{} m}{%
  \ifvmode\begin{center}\sptincludegraphics[#1]{#2}\end{center}%
  \else\sptincludegraphics[#1]{#2}\fi}
% keep figures where they are referenced
\usepackage{float}
\floatplacement{figure}{H}
% running header with the document title
\usepackage{fancyhdr}
\pagestyle{fancy}
\fancyhf{}
\fancyhead[L]{\small PROJECT spacetimes, Revision\_Deflation --- documentation and provenance}
\fancyhead[R]{\small\thepage}
\renewcommand{\headrulewidth}{0.4pt}
% 16x16 matrices: tighter columns so that a full matrix fits the text width
\setlength{\arraycolsep}{3.2pt}
% The documents of this project contain long tagged symbols, file paths and
% identifiers that cannot be hyphenated.  A paragraph is therefore allowed to
% have looser lines before TeX gives up and lets a line stick out of the margin.
\AtBeginDocument{\tolerance=2000 \emergencystretch=5em \hbadness=2000}
\ifLuaTeX
  \usepackage{selnolig}  % disable illegal ligatures
\fi
\IfFileExists{bookmark.sty}{\usepackage{bookmark}}{\usepackage{hyperref}}
\IfFileExists{xurl.sty}{\usepackage{xurl}}{} % add URL line breaks if available
\urlstyle{same}
\hypersetup{
  pdftitle={Revision\_Deflation --- the 16-component Dirac field and its singlet superfluid on spacetimes whose coordinates x5, x6, x7 deflate exponentially},
  pdfauthor={Prepared for Patrick Lee Nash, Ph.D.},
  colorlinks=true,
  linkcolor={blue},
  filecolor={Maroon},
  citecolor={blue},
  urlcolor={blue},
  pdfcreator={LaTeX via pandoc}}

\title{Revision\_Deflation --- the 16-component Dirac field and its
singlet superfluid on spacetimes whose coordinates x5, x6, x7 deflate
exponentially}
\usepackage{etoolbox}
\makeatletter
\providecommand{\subtitle}[1]{% add subtitle to \maketitle
  \apptocmd{\@title}{\par {\large #1 \par}}{}{}
}
\makeatother
\subtitle{PROJECT spacetimes: where to start, what was found, and what
is open}
\author{Prepared for Patrick Lee Nash, Ph.D.}
\date{November 13, 2026}

\begin{document}
\maketitle

\hypertarget{what-this-folder-is}{%
\section{What this folder is}\label{what-this-folder-is}}

The main project (the folder above this one) studies a 16-component
Dirac field on two flat eight-dimensional spacetimes, of signatures 4+4
and 7+1, and a superfluid of fermion pairs bound by a scalar field, the
Vorton-Scalar. On November 12, 2026 the author paused the main project
and instructed that every calculation be repeated under new hypotheses:

\begin{quote}
the coordinates \(x5\), \(x6\), \(x7\) deflate exponentially; Minkowski
space is replaced by a pseudo-Riemannian space with a diagonal metric
whose components along these coordinates are
\(g_{55}=-\mathrm{const}^2e^{-2\,\mathrm{g55}\,x4}\), and likewise
\(g_{66}\) with the rate \(\mathrm{g66}\) and \(g_{77}\) with the rate
\(\mathrm{g77}\).
\end{quote}

Here \(\mathrm{g55}\), \(\mathrm{g66}\), \(\mathrm{g77}\) are the
author's names for three positive rates, and \(g_{55}\), \(g_{66}\),
\(g_{77}\) are components of the metric. This folder contains the
repetition: programs, data, figures, documents and two manuscripts.
Nothing in the folder above was changed, except one line of a build
script (recorded in the provenance P11 of the main project).

\hypertarget{where-to-start}{%
\section{Where to start}\label{where-to-start}}

\begin{longtable}[]{@{}
  >{\raggedright\arraybackslash}p{(\columnwidth - 2\tabcolsep) * \real{0.2353}}
  >{\raggedright\arraybackslash}p{(\columnwidth - 2\tabcolsep) * \real{0.7647}}@{}}
\toprule\noalign{}
\begin{minipage}[b]{\linewidth}\raggedright
if you want
\end{minipage} & \begin{minipage}[b]{\linewidth}\raggedright
read
\end{minipage} \\
\midrule\noalign{}
\endhead
\bottomrule\noalign{}
\endlastfoot
the two papers &
\texttt{manuscript/16\ Component\ Dirac\ Fermion\ Condensate\ Singlet\ Superfluid\ on\ 4+4\ Deflating\ Spacetime.pdf}
and \texttt{...\ on\ 7+1\ Deflating\ Spacetime.pdf} \\
the geometry and the Dirac equations, derived step by step &
\texttt{dirac\ matrices\ deflating.pdf} (DMD) \\
the superfluid on the deflating spacetimes &
\texttt{supplement/RS1\_vorton\_scalar\_and\_superfluid\_deflating.pdf} \\
the effective action along the deflation &
\texttt{supplement/RS4\_effective\_action\_along\_the\_clock.pdf} \\
the Higgs hypothesis &
\texttt{supplement/RS2\_higgs\_hypothesis\_deflating.pdf} \\
cosmology, dark matter, the Hubble tension &
\texttt{supplement/RS3\_cosmology\_deflating.pdf} \\
the microwave background and the growth of structure &
\texttt{supplement/RS5\_microwave\_background\_deflating.pdf} \\
what was assumed, what mirrors what, what is open &
\texttt{provenance/RP0\_hypotheses\_and\_status.pdf} \\
the proof that no symbol has two meanings &
\texttt{provenance/RP1\_master\_symbols\_deflation.pdf} \\
how to run everything & \texttt{provenance/RP8\_code\_deflation.pdf} \\
every error that was made &
\texttt{provenance/RP9\_error\_log\_deflation.pdf} \\
\end{longtable}

Every document exists as \texttt{.md}, \texttt{.tex} and \texttt{.pdf}.

\hypertarget{what-was-found}{%
\section{What was found}\label{what-was-found}}

\textbf{Geometry} (DMD, Part 1). The deflating spacetimes are
homogeneous, with 14, 15 or 20 symmetries (three different, two equal,
three equal rates) instead of the 36 of flat spacetime. They are not
vacuum solutions of Einstein's equations. \emph{With three equal rates
they are exact vacuum solutions of Einstein--Gauss--Bonnet gravity}, for
a cosmological constant \(3\,\mathrm{g55}^2\) and a Gauss--Bonnet
coupling \(-1/(4\,\mathrm{g55}^2)\).

\textbf{The Dirac field} (DMD, Parts 2 to 4). A field that does not
depend on \(x5,x6,x7\) obeys, after a rescaling by the square root of
the deflating volume, exactly the flat equations. A field that does
depend on them behaves in opposite ways on the two spacetimes. On
\textbf{7+1} the deflation creates particles, with a Fermi--Dirac
spectrum of the temperature \(\mathrm{rate}/2\pi\). On \textbf{4+4}
every admissible mode that depends on a deflating time is a standing
wave that is extinguished as the exponential of an exponential:
\emph{the deflation removes the dependence on the additional times,
which the main project had to impose by hand.}

\textbf{The superfluid} (RS1, RS4). The model of the main project holds
unchanged, with a coupling that grows:
\(\alpha_{VS}\propto e^{\tau_d}\), where \(\tau_d=\mathrm{gsum}\,x4\) is
the number of e-foldings of the deflating volume; the masses are
constant. Every threshold of the main project is therefore passed at a
definite time. The gap grows much faster than the coupling,
\(d\ln(\mathrm{gap})/d\tau_d=1/\lambda\). Two fermions become bound at a
scattering resonance. The paired state can follow the deflation only
once the gap exceeds a few times \(\mathrm{gsum}/\lambda\); before that
it lags and breaks pairs. A vortex line carries four fermion zero modes,
two for each direction.

\textbf{The Higgs hypothesis} (RS2). The internal symmetry that gives
the pair condensate the pattern of the Higgs vacuum survives exactly,
but on the deflating 7+1 spacetime it is no longer a symmetry of the
spacetime. And if the condensate were the Higgs field, the electroweak
scale would follow the growing coupling; atomic clocks then limit the
rate of deflation to about \(10^{-6}\) of the Hubble constant. \emph{The
Higgs hypothesis and a deflation of any consequence today exclude each
other.}

\textbf{Cosmology} (RS3). If gravity is the eight-dimensional Einstein
theory, Newton's constant grows with the deflation, lunar laser ranging
limits the rate to \(3\times10^{-3}\) of the Hubble constant, and the
Hubble constant inferred from the microwave background \emph{falls}. The
Hubble tension is not relieved. For dark matter the deflation adds an
epoch of resonant self-interaction; the conclusions of the main project
about the masses are unchanged.

\textbf{The microwave background and the growth of structure} (RS5). The
spectra need one assumption more: that the scale factors of \(x5,x6,x7\)
are not perturbed when ordinary space is. Even then the Einstein
equations of ordinary space cannot all hold with matter alone. Whatever
keeps the three rates constant must carry a tension and perturbations of
it, and the hypotheses do not say which. \emph{What does not depend on
that choice:} the spectra at the multipoles \(\ell\ge30\). Mock data of
the quality of the Planck satellite would bound the rate to
\(0.0385\,H_0\), between the bounds of lunar laser ranging and of
nucleosynthesis; in every fit the Hubble constant falls, by 84 km/s/Mpc
per unit of \(\mathrm{gsum}/H_0\); and the rate cannot help the dark
radiation of the main project to raise the Hubble constant. \emph{What
does depend on it:} the growth of structure and the largest angular
scales, which the revision brackets with two choices.

\hypertarget{what-is-open}{%
\section{What is open}\label{what-is-open}}

\begin{enumerate}
\def\labelenumi{\arabic{enumi}.}
\tightlist
\item
  What keeps the three rates constant. The revision calls it the agent
  and assumes nothing about its nature. The spectra of the microwave
  background were computed under an assumption about it that the author
  is asked to confirm or replace (RP0, Sec. 5); the growth of structure
  remains undetermined until it is specified.
\item
  The stability of the Einstein--Gauss--Bonnet solution, and the source
  of the deflation for unequal rates.
\item
  The back-reaction of the created particles (7+1), and a principle that
  excludes the growing solutions (4+4).
\item
  The end of every history: the couplings grow without bound, and the
  model fails when they become strong.
\end{enumerate}

The stages of the main project that were paused are listed in its
provenance P11 and remain paused.

\hypertarget{numbers}{%
\section{Numbers}\label{numbers}}

The Rust program \texttt{deflation} performs 238 checks and has 49 unit
tests; an independent Python script performs 43 of 43 checks; the
registry of symbols 6 of 6. All pass. A further script compares every
number that is quoted in the running text of the documents and of the
manuscripts with the numbers that the programs state (974 numbers in 29
documents, 0 not traced). The folder contains 19 figures and the data
files from which every table and figure is made.

\hypertarget{how-to-run-it}{%
\section{How to run it}\label{how-to-run-it}}

\begin{Shaded}
\begin{Highlighting}[]
\BuiltInTok{cd} \StringTok{"/Users/nsh/Developer/github/spacetimes/Revision\_Deflation"}
\ExtensionTok{python3}\NormalTok{ code/scripts/vendor\_eb.py }\AttributeTok{{-}{-}check}
\ExtensionTok{cargo}\NormalTok{ build }\AttributeTok{{-}{-}release} \AttributeTok{{-}{-}manifest{-}path}\NormalTok{ code/deflation{-}rs/Cargo.toml}
\ExtensionTok{cargo}\NormalTok{ test  }\AttributeTok{{-}{-}release} \AttributeTok{{-}{-}manifest{-}path}\NormalTok{ code/deflation{-}rs/Cargo.toml}
\CommentTok{\# the first nine commands: about three and a half minutes together}
\ControlFlowTok{for}\NormalTok{ c }\KeywordTok{in}\NormalTok{ matrices modes history quench cosmology symmetry channels matter paired}\KeywordTok{;} \ControlFlowTok{do}
  \ExtensionTok{./code/deflation{-}rs/target/release/deflation} \VariableTok{$c}
\ControlFlowTok{done}
\CommentTok{\# the Einstein{-}Boltzmann calculations of RS5: about an hour and ten minutes together}
\ControlFlowTok{for}\NormalTok{ c }\KeywordTok{in}\NormalTok{ background spectrum forecast families superfluid}\KeywordTok{;} \ControlFlowTok{do}
  \ExtensionTok{./code/deflation{-}rs/target/release/deflation} \VariableTok{$c}
\ControlFlowTok{done}
\ExtensionTok{python3}\NormalTok{ code/scripts/verify\_deflating\_independent.py}
\end{Highlighting}
\end{Shaded}

The provenance RP8 explains every step, and the remaining ones (symbols,
tables, figures, documents).

\end{document}
